\section*{第 \ref{ch:b1:euclid} 章——欧几里得空间}

\begin{solution}{exo:b1:euclid:1}
$\langle u, v\rangle = 2 - 4 + 2 = 0$；$\norm u = \norm v = 3$。
两向量正交：夹角为 $\frac\pi2$。柯西–施瓦茨：
$\abs 0 \leq 9$，显然成立。
\end{solution}

\begin{solution}{exo:b1:euclid:2}
将两个平方范数展开，利用
\cref{thm:b1:euclid:cs}中的恒等式，再相加：交叉项相消。

无穷范数：取 $x = (1, 0)$，$y = (0, 1)$。则 $\norm{x + y}_\infty
= \norm{x - y}_\infty = 1$，而恒等式要求 $1 + 1 =
2(1) + 2(1) = 4$，不成立。违反平行四边形
恒等式的范数不来自任何\omterm{def:b1:euclid:def}{内积}。
\end{solution}

\begin{solution}{exo:b1:euclid:3}
$e_1 = \frac{1}{\sqrt 2}(1,1,0)$.

$w_2 = (1,0,1) - \langle (1,0,1), e_1\rangle e_1 = (1,0,1) -
\frac{1}{2}(1,1,0) = \bigl(\tfrac12, -\tfrac12, 1\bigr)$;
$\norm{w_2} = \sqrt{\tfrac32}$，于是 $e_2 =
\frac{1}{\sqrt 6}(1, -1, 2)$.

$w_3 = (0,1,1) - \langle\cdot, e_1\rangle e_1 - \langle\cdot,
e_2\rangle e_2 = (0,1,1) - \frac12 (1,1,0) - \frac16(1,-1,2) =
\bigl(-\tfrac23, \tfrac23, \tfrac23\bigr)$；单位化，$e_3 =
\frac{1}{\sqrt 3}(-1, 1, 1)$.

（验证：两两\omterm{def:b1:euclid:def}{内积}为零；每个向量的范数均为 $1$。）
\end{solution}

\begin{solution}{exo:b1:euclid:4}
$F^\perp = \{(x,y,z) : x + y + z = 0\}$，基为 $\bigl((1,-1,0),
(1,0,-1)\bigr)$。取 $e = \frac{1}{\sqrt3}(1,1,1)$，则 $p_F(x) =
\langle x, e\rangle e = \frac{x_1 + x_2 + x_3}{3}\,(1,1,1)$，于是
\[
\operatorname{Mat}(p_F) = \frac13
\begin{pmatrix} 1&1&1\\ 1&1&1\\ 1&1&1\end{pmatrix}.
\]
当 $x = (1,2,3)$ 时，$p_F(x) = (2,2,2)$，且 $d(x, F) = \norm{(1,2,3)
- (2,2,2)} = \norm{(-1,0,1)} = \sqrt 2$.
\end{solution}

\begin{solution}{exo:b1:euclid:5}
设 $p = \alpha(1,0,1) + \beta(0,1,1) = (\alpha, \beta, \alpha +
\beta)$。$x - p$ 与各生成元正交：
\[
\langle x - p, (1,0,1)\rangle = (1 - \alpha) + (4 - \alpha - \beta)
= 0,
\qquad
\langle x - p, (0,1,1)\rangle = (1 - \beta) + (4 - \alpha - \beta)
= 0,
\]
即 $2\alpha + \beta = 5$ 且 $\alpha + 2\beta = 5$：$\alpha =
\beta = \frac53$。于是 $p_F(x) = \bigl(\frac53, \frac53,
\frac{10}{3}\bigr)$，且
\[
d(x, F) = \norm{x - p} =
\norm{\bigl(-\tfrac23, -\tfrac23, \tfrac23\bigr)} =
\frac{2}{\sqrt 3} .
\]
格拉姆–施密特正交化并非必要，因为\omterm{def:b1:linmaps:projection}{投影}的定义性质
——$x - p \perp F$——本身是一个\omterm{def:b1:det:system}{线性方程组}
（“正规方程”），关于\emph{任一}生成
族中的系数。
\end{solution}

\begin{solution}{exo:b1:euclid:6}
柯西–施瓦茨不等式，取 $g = 1$：
\[
\Bigl(\int_0^1 f \cdot 1\Bigr)^2 \leq \int_0^1 f^2 \cdot \int_0^1
1^2 = \int_0^1 f^2 ,
\]
等号成立当且仅当 $f$ 与 $1$ 成比例，即为常数。离散
形式：在 $\R^n$ 中，取 $x = (a_1, \dots, a_n)$，$y = (1, \dots,
1)$:
$\bigl(\sum a_i\bigr)^2 \leq \norm x^2\, \norm y^2 = n \sum a_i^2$,
等号成立当且仅当所有 $a_i$ 相等。
\end{solution}

\begin{solution}{exo:b1:euclid:7}
由\cref{thm:b1:euclid:projection}，$E = F \oplus F^\perp$，于是
维数相加：$\dim F^\perp = \dim E - \dim F$。包含关系 $F
\subseteq (F^\perp)^\perp$ 显然成立（$F$ 中的向量
与 $F^\perp$ 中的一切向量正交）。维数：
$\dim(F^\perp)^\perp = \dim E - \dim F^\perp = \dim F$；
\omterm{def:b1:findim:def}{有限维}时，维数相等的包含关系即为相等
（见\cref{thm:b1:findim:subspaces}）。
\end{solution}

\begin{solution}{exo:b1:euclid:8}
$A$：各列为单位向量，行列式为 $+1$，故为旋转，其中 $\cos\theta =
\frac{1}{\sqrt2}$, $\sin\theta = \frac{1}{\sqrt2}$: $\theta =
\frac\pi4$，因此 $A^8 = R_{8\pi/4} = R_{2\pi} = I$。

$B$：行列式为 $\frac{1}{25}(-9 - 16) = -1$，是反射
$S_\theta$，其中 $\cos\theta = \frac35$，$\sin\theta = \frac45$；其
对称轴的倾角为 $\frac\theta2$（相对于 $x$ 轴），即
以 $\bigl(\cos\frac\theta2, \sin\frac\theta2\bigr)$ 为方向向量的直线
——具体而言，对称轴为
$\operatorname{Vect}\bigl((2, 1)\bigr)$，因为 $B(2,1)^{\mathsf T} =
\frac15(6+4, 8-3)^{\mathsf T} = (2,1)^{\mathsf T}$。作为反射，
$B^2 = I$.
\end{solution}

\begin{solution}{exo:b1:euclid:9}
所求量为 $\norm{f - (a + bX)}^2$，它在 $C(\intcc{0}{1})$ 中赋以
\omterm{thm:b1:integration:def}{积分}\omterm{def:b1:euclid:def}{内积}：恰在
\omterm{thm:b1:euclid:projection}{正交投影} $a + bX = p_F(f) = X - \frac16$ 处取得
（见\cref{ex:b1:euclid:bestapprox}）。最小值为
\[
\norm{f - p_F(f)}^2 = \int_0^1 \Bigl(x^2 - x + \frac16\Bigr)^{\!2}
\dd x .
\]
展开：$\int_0^1 (x^2 - x + \frac16)^2 = \int_0^1 \bigl(x^4 - 2x^3
+ \frac{4}{3}x^2 - \frac13 x + \frac{1}{36}\bigr)\dd x = \frac15 -
\frac12 + \frac49 - \frac16 + \frac{1}{36} = \frac{1}{180}$.
于是最小值为 $\dfrac{1}{180}$。
\end{solution}

\begin{solution}{exo:b1:euclid:10}
设 $y \in \ker(u - \mathrm{id})$（即 $u(y) = y$），且 $x \in E$。
于是利用\omterm{def:b1:euclid:def}{内积}的保持性（$\langle u(a),
u(b)\rangle = \langle a, b\rangle$）：
\[
\langle x - u(x),\, y\rangle
= \langle x, y\rangle - \langle u(x), y\rangle
= \langle x, y\rangle - \langle u(x), u(y)\rangle
= \langle x, y\rangle - \langle x, y\rangle = 0 :
\]
这说明 $\operatorname{im}(u - \mathrm{id})$ 中的每个向量都与
每个不动向量正交。

因此 $\operatorname{im}(u - \mathrm{id}) \subseteq \ker(u -
\mathrm{id})^{\perp}$。由秩--零度定理及
\cref{exo:b1:euclid:7}，二者维数均为 $\dim E - \dim\ker(u -
\mathrm{id})$，因此二者相等。于是 $E = \ker(u - \mathrm{id})
\oplus \ker(u - \mathrm{id})^\perp = \ker(u - \mathrm{id}) \oplus
\operatorname{im}(u - \mathrm{id})$。
\end{solution}

\begin{solution}{exo:b1:euclid:11}
\begin{enumerate}
  \item 设 $Gc = 0$，其中列向量 $c = (c_1, \dots,
        c_k)^{\mathsf T}$。则
        \[
        \Bigl\lVert\sum_i c_i v_i\Bigr\rVert^2
        = \sum_{i,j} c_i c_j \langle v_i, v_j\rangle
        = c^{\mathsf T} G\, c = 0,
        \]
        于是 $\sum_i c_i v_i = 0$。若该族\omterm{def:b1:vspaces:free}{线性无关}，则
        必有 $c = 0$，即 $G$ 可逆。反之，
        若存在非平凡关系 $\sum_j c_j v_j = 0$，则分别与各 $v_i$ 作\omterm{def:b1:euclid:def}{内积}，
        得非平凡关系
        $Gc = 0$，故 $G$ 奇异。\omterm{def:b1:vspaces:free}{线性无关} $\iff \det G \neq 0$。
  \item $\langle X^{i-1}, X^{j-1}\rangle = \int_0^1 x^{i+j-2}\dd
        x = \frac{1}{i+j-1}$：$(1, X, X^2)$ 的格拉姆矩阵
        恰为\omterm{pb:b1:det:1}{希尔伯特矩阵} $H_3$，其行列式
        $\frac1{2160}$ 由\cref{ch:b1:det}的周末问题算出：
        非零，故单项式\omterm{def:b1:vspaces:free}{线性无关}——
        与预期相符，但现在由一个数值加以证实。
\end{enumerate}
\end{solution}

\begin{solution}{exo:b1:euclid:12}
\begin{enumerate}
  \item $A^2 = A A = A^{\mathsf T} A = I$，故 $u^2 =
        \mathrm{id}$：$u$ 是对称，且
        由\cref{thm:b1:linmaps:projchar}~(2)可得 $E = \ker(u -
        \mathrm{id}) \oplus \ker(u + \mathrm{id})$。
  \item 设 $u(x) = x$，$u(y) = -y$。由于 $u$ 保持
        \omterm{def:b1:euclid:def}{内积}，
        \[
        \langle x, y\rangle = \langle u(x), u(y)\rangle
        = \langle x, -y\rangle = -\langle x, y\rangle ,
        \]
        故 $\langle x, y\rangle = 0$，两个特征\omterm{def:b1:vspaces:subspace}{子空间}
        正交，$\ker(u + \mathrm{id}) = F^\perp$，其中 $F =
        \ker(u - \mathrm{id})$，且 $u$ 是
        关于 $F$ 的正交对称。
  \item $R_\theta$ 对称当且仅当 $-\sin\theta = \sin\theta$，
        即 $\theta \in \{0, \pi\}$：\omterm{def:b1:logic:map}{映射} $\pm\mathrm{id}$
        （恒等\omterm{def:b1:logic:map}{映射}与中心对称）。每个 $S_\theta$ 都是
        对称，即关于直线的对称。这些恰为
        平面上的正交对称，$F$ 分别取
        整个平面、$\{0\}$ 或反射轴。
\end{enumerate}
\end{solution}

\begin{solution}{pb:b1:euclid:1}
\textbf{1.} $\norm{t_a(x) - t_a(y)} = \norm{x - y}$；线性\omterm{def:b1:euclid:isometry}{等距变换}
保持范数，从而保持距离（$\norm{g(x) - g(y)} =
\norm{g(x - y)} = \norm{x-y}$）；保距\omterm{def:b1:logic:map}{映射}的复合仍保持距离。
每个\omterm{def:b1:euclid:isometry}{等距变换}都是\omterm{def:b1:logic:inj}{单射}（不同点之间保持正距离），并
由下面的分类知它为\omterm{def:b1:logic:inj}{双射}；恒等\omterm{def:b1:logic:map}{映射}与
逆\omterm{def:b1:logic:map}{映射}均为\omterm{def:b1:euclid:isometry}{等距变换}，
从而构成一个群。

\textbf{2.} $\norm{f(x)} = \norm{f(x) - f(0)} = \norm{x - 0} =
\norm x$。由极化恒等式：
\[
\langle f(x), f(y)\rangle
= \frac{\norm{f(x)}^2 + \norm{f(y)}^2 - \norm{f(x) -
f(y)}^2}{2}
= \frac{\norm x^2 + \norm y^2 - \norm{x - y}^2}{2}
= \langle x, y\rangle .
\]

\textbf{3.} 利用第 2 小题展开（$f$ 的每个乘积
等于其自变量的乘积）：
\begin{align*}
\norm{f(x+y) - f(x) - f(y)}^2
&= \norm{x+y}^2 + \norm x^2 + \norm y^2\\
&\quad - 2\langle x+y, x\rangle - 2\langle x+y, y\rangle
+ 2\langle x, y\rangle = 0 ,
\end{align*}
直接验算可知上式为零；类似地，$\norm{f(\lambda x) -
\lambda f(x)}^2 = \norm{\lambda x}^2 - 2\lambda\langle \lambda
x, x\rangle + \lambda^2\norm x^2 = 0$。于是 $f(x + y) = f(x) +
f(y)$，$f(\lambda x) = \lambda f(x)$，故 $f$ 线性，且
保持范数，即 $f \in O(\R^2)$。

\textbf{4.} 令 $a = f(0)$，$g = t_{-a}\circ f$，它是固定
$0$ 的\omterm{def:b1:euclid:isometry}{等距变换}，由第 3 小题知为线性\omterm{def:b1:euclid:isometry}{等距变换}，且 $f = t_a
\circ g$。唯一性：由 $t_a \circ g = t_{a'}\circ g'$ 在
$0$ 处取值得 $a = a'$，进而 $g = g'$。

\textbf{5.} $(t_a\circ g)\circ(t_{a'}\circ g') = t_{a +
g(a')}\circ(g\circ g')$，因为 $g\,t_{a'} = t_{g(a')}\,g$ 对
\omterm{def:b1:linmaps:def}{线性映射} $g$ 成立。于是复合的线性部分为 $g \circ g'$，
且 $\det(gg') = \det g \det g'$，于是
保向 $\circ$ 保向 $=$ 反向 $\circ$ 反向 $=$ 保向，
保向 $\circ$ 反向 $=$ 反向——即 $\pm1$ 的符号法则。

\textbf{6.} $z \mapsto \eu^{\iu\theta}z$ 将 $1$ 映到
$(\cos\theta, \sin\theta)$，而 $\iu$ 映到
$\iu\eu^{\iu\theta} = (-\sin\theta, \cos\theta)$，即
$R_\theta$ 的列。又 $z \mapsto \eu^{\iu\theta}\conj z$ 将 $1$ 映到
$(\cos\theta, \sin\theta)$，而 $\iu$ 映到 $-\iu\eu^{\iu\theta} =
(\sin\theta, -\cos\theta)$，即 $S_\theta$ 的列。

\textbf{7.} 当 $a = 1$ 时：$f = t_b$。当 $a \neq 1$ 时，不动点
方程 $az_0 + b = z_0$ 有唯一解 $z_0 = b/(1 -
a)$，于是
\[
f(z) - z_0 = az + b - (az_0 + b) = a(z - z_0):
\]
在以 $z_0$ 为中心的\omterm{prop:b1:vspaces:coordinates}{坐标}系中，$f$ 是乘以 $a
= \eu^{\iu\arg a}$，即中心为 $z_0$ 的
旋转，旋转角为 $\arg a$。

\textbf{8.} $f(f(z)) = a\conj{(a\conj z + b)} + b = a\conj a\,z
+ a\conj b + b = z + v$，故 $f^2 = t_v$。可交换性：$f(z + v) =
a\conj z + a\conj v + b$，且 $a\conj v = a\conj{(a\conj b + b)}
= b + a\conj b = v$，从而 $f\circ t_v = t_v\circ f$。

\emph{情形 $v = 0$}：$f^2 = \mathrm{id}$。对任意 $z$，
中点 $m = \frac{z + f(z)}2$ 满足（\omterm{def:b1:euclid:isometry}{等距变换}是仿射的，
见第一部分）$f(m) = \frac{f(z) + f^2(z)}{2} = m$：不动点
存在。以不动点为新原点作平移共轭，
$f$ 化为线性反向\omterm{def:b1:euclid:isometry}{等距变换}，即某个 $S_\theta$
（见\cref{thm:b1:euclid:plane}），即关于一条直线的对称。

\emph{情形 $v \neq 0$}：$r = t_{-v/2}\circ f$ 为反向变换，且
\[
r^2 = t_{-v/2}\,f\,t_{-v/2}\,f = t_{-v/2}\,t_{-v/2}\,f^2
= t_{-v}\,t_v = \mathrm{id}
\]
（利用可交换性），故 $r$ 为一个对称，且 $f =
t_{v/2}\circ r$。其轴平行于 $v$：$r$ 与
$t_v$ 可交换（$f$ 与 $t_{v/2}$ 均如此），故 $t_v$ 将轴（
$r$ 的不动直线）映为自身，迫使 $v$ 沿轴方向。由此
$f$ 是滑移向量为 $v/2$ 的滑移对称。至此，每个
\omterm{def:b1:euclid:isometry}{等距变换}都已分类：平移或旋转（保向），
对称或滑移对称（反向）。

\textbf{9.} 旋转 $f(z) = az + b$，$g(z) = a'z + b'$，其中
$a = \eu^{\iu\alpha}$，$a' = \eu^{\iu\beta}$，则 $g\circ f(z) =
a'az + (a'b + b')$ 的线性系数为
$\eu^{\iu(\alpha+\beta)}$，是角为 $\alpha + \beta$ 的旋转，
由第 7 小题知，或当 $\eu^{\iu(\alpha+\beta)} =
1$ 时为平移。对称 $r_i(z) = a_i\conj z + b_i$，其中 $a_i =
\eu^{\iu\theta_i}$（轴的倾角为 $\theta_i/2$）：
\[
r_2\circ r_1(z) = a_2\conj{a_1}\,z + (a_2\conj{b_1} + b_2),
\]
保向，旋转角为 $\theta_2 - \theta_1 = 2\bigl(\tfrac{\theta_2}2
- \tfrac{\theta_1}2\bigr)$，即两轴夹角的两倍；
两轴平行（$\theta_1 = \theta_2$）时给出平移。

\textbf{10.} 中心为 $c$、旋转角为 $\theta$ 的旋转，是
两条过 $c$ 的直线上两个对称之积，两直线夹角为
$\theta/2$（由第 9 小题逆推）；平移 $t_v$ 是
两条与 $v$ \omterm{def:b1:euclid:orthogonal}{正交的}平行直线上两个对称之积，
两直线相距 $\norm v/2$；对称即一个对称；
滑移对称是一个对称与一个平移的复合，
因此需要三个。最多即为三个。

\textbf{11.} $a = \iu$, $b = 1 + \iu$: $v = a\conj b + b =
\iu(1 - \iu) + 1 + \iu = (1 + \iu) + (1 + \iu) = 2 + 2\iu \neq
0$，故为滑移向量为 $v/2 = 1 + \iu$ 的滑移对称。
记对称 $r = t_{-(1+\iu)}\circ f$，则 $r(z) = \iu\conj z$，
其轴是倾角为 $\frac12\arg\iu = \frac\pi4$ 的直线：
即直线 $y = x$。故 $f$ 是轴为 $y =
x$、滑移向量为 $(1, 1)$ 的滑移对称。

\textbf{12.} $f(z) = \iu z$，且 $g(z) = 1 + \iu(z - 1) = \iu z
+ 1 - \iu$，于是
\[
g\circ f(z) = \iu(\iu z) + 1 - \iu = -z + 1 - \iu :
\]
线性系数为 $-1 = \eu^{\iu\pi}$，是角为
$\pi$ 的旋转（即半圈旋转），中心为 $z_0 = \frac{1 - \iu}{1 - (-1)}
= \frac{1 - \iu}{2}$。

\textbf{13.} $r_2\circ r_1(z) = \iu\,\conj{\conj z} = \iu z$:
即旋转角为 $\frac\pi2$、中心为 $0$ 的旋转。两轴（实轴，
倾角为 $0$；直线 $y = x$，倾角为 $\frac\pi4$）相交于
夹角为 $\frac\pi4$，其两倍为 $\frac\pi2$，与第 9 小题的结论
相符。

\textbf{14.} 设 $f = t_a\circ g$，$g$ 为\omterm{def:b1:linmaps:def}{线性映射}，且 $\sum_i
\lambda_i = 1$：
\[
f\Bigl(\sum_i \lambda_i x_i\Bigr)
= a + \sum_i \lambda_i\,g(x_i)
= \sum_i \lambda_i\bigl(a + g(x_i)\bigr)
= \sum_i \lambda_i f(x_i) .
\]

\textbf{15.} 对 $h \in G$，利用第 14 小题（系数
$\frac1{\abs G}$ 之和为 $1$）：
\[
h(c) = \frac{1}{\abs G}\sum_{g \in G} h\bigl(g(x_0)\bigr)
= \frac{1}{\abs G}\sum_{g' \in G} g'(x_0) = c ,
\]
因为 $g \mapsto hg$ 是 $G$ 到自身的\omterm{def:b1:logic:inj}{双射}。

\textbf{16.} 共轭变换 $t_{-c}\,g\,t_{c}$（$g \in G$）构成
一个固定 $0$ 的\omterm{def:b1:euclid:isometry}{等距变换}群，因而为\emph{线性}
\omterm{def:b1:euclid:isometry}{等距变换}（由第 3 小题）：$O(\R^2)$ 的有限\omterm{def:b1:structures:subgroup}{子群}。若
存在 $s \in G$ 为反向的，则\omterm{def:b1:logic:map}{映射} $g \mapsto sg$ 把 $G^+$
\omterm{def:b1:logic:inj}{单射}地映入 $G\setminus G^+$，而 $h \mapsto s^{-1}h$ 为其
逆\omterm{def:b1:logic:map}{映射}（由第 5 小题的符号法则）：$\abs G = 2\,\abs{G^+}$；
否则 $G = G^+$。

\textbf{17.} 若 $G^{+} = \{\mathrm{id}\}$，则其为 $C_1$。
否则将其元素记作 $R_\theta$，$\theta \in
\intco{0}{2\pi}$，设 $\theta_0$ 为出现的最小正
角。对 $R_\theta \in G^{+}$，作带余除法
$\theta = k\theta_0 + \rho$，其中 $0 \leq \rho < \theta_0$，可得
$R_\rho = R_\theta R_{\theta_0}^{-k} \in G^{+}$，由极小性得 $\rho = 0$，
于是 $G^{+} = \langle R_{\theta_0}\rangle$。
将 $2\pi$ 除以 $\theta_0$，用同样方法得 $2\pi = n\theta_0$：
$G^{+} \cong C_n$，由角为
$\frac{2\pi}n$ 的旋转生成。

\textbf{18.} 由第 16 小题的计数，$G = G^{+} \cup sG^{+}$。
取 $s(z) = a\conj z$，$r(z) = \omega z$
（\omterm{def:b1:linmaps:forms}{线性型}，由第 16 小题可知；$\abs a = \abs\omega = 1$）：
\[
s r s(z) = a\,\conj{\omega\,a\conj z}
= a\conj\omega\conj a\,z = \conj\omega\,z = r^{-1}(z) .
\]
每个元素 $sr^k$ 属于 $sG^{+}$，它是反向的，并满足
$(sr^k)^2 = (sr^ks)\,r^k = r^{-k}r^k = \mathrm{id}$：一个对合的
反向\omterm{def:b1:euclid:isometry}{等距变换}，固定 $c$，即反射。于是
$G$ 由 $n$ 个旋转与 $n$ 个反射组成，其关系为
$r^n = s^2 = \mathrm{id}$，$srs = r^{-1}$：这是
\omterm{pb:b1:euclid:1}{二面体群} $D_n$。

\textbf{19.} 综合可得：平面\omterm{def:b1:euclid:isometry}{等距变换}的有限群，
固定某点 $c$（第 15 小题），并化为
$O(2)$ 的有限\omterm{def:b1:structures:subgroup}{子群}（第 16 小题），若仅含
旋转（第 17 小题）时为 $C_n$，否则为 $D_n$（第 18 小题）：
\emph{莱昂纳多定理}。

\textbf{20.} 平移 $t_v \neq \mathrm{id}$ 的幂为
$t_{kv}$ 互不相同（$kv$ 两两不等，前提是 $v \neq 0$）：
阶无限。\omterm{pb:b1:euclid:1}{滑移反射} $f$ 满足 $f^2 = t_v$，其中 $v
\neq 0$，同样阶无限。二者都不可能属于有限群，
——与第 15--19 小题相符，那些小题给出的只有
旋转与反射。

\textbf{21.} 这 $2n$ 个\omterm{def:b1:logic:map}{映射}保持顶点集：$\omega^k
\omega^j = \omega^{k+j}$ 及 $\omega^k\conj{\omega^j} =
\omega^{k-j}$；它们是\omterm{def:b1:euclid:isometry}{等距变换}，故保持多边形
（线段映为线段）。反之，对称变换保持顶点的
重心 $0$（第 14 小题），因而是线性的，
并\omterm{def:b1:counting:objects}{置换}顶点（它们是多边形中
到 $0$ 距离最大的点）。把
顶点 $1$ 送到 $\omega^k$ 的线性\omterm{def:b1:euclid:isometry}{等距变换}为 $z \mapsto \omega^k z$，若为正向，
$z \mapsto \omega^k\conj z$，若为反向（其矩阵由
一列及符号唯一确定）：至多 $2n$
个对称变换，故恰为上面的 $D_n$。

\textbf{22.} $C_1$：不等边三角形（无非平凡
对称）。$D_1$：等腰非等边三角形（一个
反射）。$D_2$：非正方形的矩形（恒等变换、
绕中心的半转、两个轴反射）。$C_3$：一个
三旋臂图形——三条全等的弯臂，以 $120$ 度
为间隔相连；弯折消去了全部反射，只保留
$3$ 阶旋转。

\textbf{23.} 取 $\omega = \iu$，则旋转为 $z \mapsto z$，
$\iu z$、$-z$、$-\iu z$（旋转角分别为 $0, \frac\pi2, \pi,
\frac{3\pi}2$），以及下列反射
\[
z \mapsto \conj z \ (\text{对称轴 } y = 0), \quad
\iu\conj z \ (y = x), \quad
-\conj z \ (x = 0), \quad
-\iu\conj z \ (y = -x):
\]
正方形的两条对角线与两条中线——\omterm{pb:b1:euclid:1}{二面体群}
$D_4$，阶为 $8$。

\textbf{24.} 记 $a = \eu^{\iu\theta} \neq 1$，$f(z) = az +
c_1(1 - a)$，$g(z) = az + c_2(1 - a)$。则
\[
f\circ g(z) - g\circ f(z)
= (1-a)\bigl(ac_2 + c_1 - ac_1 - c_2\bigr)
= -(1 - a)^2(c_2 - c_1) \neq 0 ,
\]
故 $f g \neq g f$。两个复合\omterm{def:b1:logic:map}{映射}的线性系数均为
$a^2$，故 $(fg)\circ(gf)^{-1}$ 的线性系数为 $1$，它是
一个平移，平移量即 $fg(z) - gf(z)$，也就是上面的非零常数。
含 $f$ 与 $g$ 的群必含此平移及
其全部幂，故为无限群。两个旋转中心，对有限群而言
已多了一个——这恰是莱昂纳多定理的
几何核心。

\textbf{25.} (i) 第 3--4 小题表明，每个保距\omterm{def:b1:logic:map}{映射}
都是仿射的，且线性部分正交；第 15 小题
断定每个有限群都固定某一点。两个结果都把度量几何
在原点处化为线性代数。(ii) 完整清单是
平移、旋转、反射、\omterm{pb:b1:euclid:1}{滑
移反射}；线性部分的行列式区分正向
与反向，而不动点的存在与否区分了
每种奇偶性中的两类。(iii) 由第 9 小题，两个反射
可复合得到任一旋转或平移，因此
反射生成整个群——至多三个即可，
对任一\omterm{def:b1:euclid:isometry}{等距变换}皆然。(iv) 在有限情形下，仅
留存循环群与\omterm{pb:b1:euclid:1}{二面体群}，这正是旋转对称
花饰恰有两类（带或不带
镜面对称轴）。第二部分证明了\emph{平面\omterm{def:b1:euclid:isometry}{等距变换}的
分类}；第四部分证明了\emph{莱昂纳多定理}。
\end{solution}
